% What my project actually is -- a plain explanation of RHOB for a reader who
% has not seen it. Single file, drops straight into Overleaf, compiles with
% pdflatex on a base TeX Live install.
%
% Every number here comes from the repository: README.md, docs/THREAT_MODEL.md,
% leaderboard/v5_replicated.json, and the module docstrings under src/rhob/.

\documentclass[11pt,a4paper]{article}

\usepackage[margin=1.1in]{geometry}
\usepackage[T1]{fontenc}
\usepackage[utf8]{inputenc}
\usepackage{booktabs}
\usepackage{xcolor}
\usepackage[colorlinks=true,linkcolor=blue!45!black,urlcolor=blue!45!black]{hyperref}

% Plain-document layout: no paragraph indents, space between paragraphs.
\setlength{\parindent}{0pt}
\setlength{\parskip}{0.7em plus 0.15em minus 0.05em}

\newcommand{\code}[1]{\texttt{\small #1}}

% Unnumbered sans-serif headings, so this reads as a document and not a paper.
\newcommand{\heading}[1]{%
  \par\vspace{1.5em}%
  \phantomsection\addcontentsline{toc}{section}{#1}%
  {\sffamily\bfseries\large #1}%
  \par\vspace{0.15em}%
  \hrule height 0.4pt%
  \par\vspace{0.5em}%
}
\newcommand{\subhead}[1]{%
  \par\vspace{0.9em}%
  {\sffamily\bfseries #1}%
  \par\vspace{0.15em}%
}

\begin{document}

\thispagestyle{empty}

{\sffamily\bfseries\Large What my project actually is}

\vspace{0.4em}
{\sffamily Aarav Shah \quad\textbullet\quad August 2026 \quad\textbullet\quad
\texttt{ashah264@ucr.edu}}

\vspace{0.3em}
\hrule height 0.8pt
\vspace{1.2em}

John,

You asked me to explain the purpose of this project, the techniques in it, and why
it matters. Here it is. I have kept the words simple and defined each one the first
time I use it. Where I got something wrong I have said so in the main text instead
of hiding it at the bottom.

Your questions are listed at the end with the page where I answer each one.

\heading{The short version}
\label{h:short}

An agent is trained to make a score go up. A person writes that score in code. The
score is never exactly the thing the person wanted. So the agent can make the score
go up by finding a flaw in it instead of doing the task. That is reward hacking.

Now say you want to write a program that looks at a finished training run and says
``this agent was gaming the score'' or ``this agent did the task''. How would you
know if your program works?

You need examples where you already know the answer. That is what I built. There
are 33 small task families. Each one makes two runs of the same task. In one the
agent games the score. In the other it does the task. Both earn about the same
visible score, on purpose. Then 30 candidate programs try to tell the two apart,
and I score how often they get it right.

I published a headline number of 0.994 out of 1. Then I went back and checked my
own setup, and the number was fake. It came from a formatting rule I had written in
my own contributor guide, not from anything about reward hacking. Take that rule
away and the number drops to 0.508, which is a coin flip.

The rest of this explains each piece, and then what is left standing.

\heading{The words I was using}
\label{h:words}

You said my words could mean a lot of different things. Here is what each one means
here.

\subhead{Proxy reward}
The number the agent is actually trying to make big. It is a real function in
Python. The environment computes it. I record it once per episode.

\subhead{True reward}
A second number, also written in Python by whoever built the task. It measures what
I actually wanted the agent to do. The agent never sees it and never optimises it.
It exists so that I know which run was the honest one.

\subhead{Oracle}
Just a name for the true reward. It is not a model, not an LLM, not a person. It is
the second scoring function above. I should have said ``true reward'' from the
start. The word oracle adds nothing and I will stop using it.

\subhead{Run}
One full recorded training run, of one agent, on one variant. You asked whether
this is a policy, a trajectory, or a distribution over trajectories. It is none of
those. It is a time series over training. In the small gridworld task it is 500
episodes, saved as four arrays of length 500: the proxy reward each episode, the
true reward each episode, a state visit count each episode, and one behaviour
number each episode.

\subhead{Cheating run and honest run}
A cheating run comes from the arm of the task that games the score. An honest run
comes from the arm that does the task. The two arms together are a \emph{pair}.

\subhead{Family}
A generator that makes those pairs. There are 33. I say family rather than
environment because one family makes many pairs, at five difficulty settings and
many random layouts.

\subhead{Cell}
One family at one difficulty. There are 123 of them. A cell is the thing that gets
scored.

\subhead{Entry check}
A test a pair has to pass before it is allowed into the set at all. This is the
thing I was calling ``admission criteria''. You were right that the phrase is bad.
Entry check is what I meant, and it is the part of the project I care about most.

\heading{Do the two arms learn? Mostly no}
\label{h:learn}

You asked whether the learning algorithms are the same, whether the reward
functions are the same, and whether these can be seen as two seeds optimising the
same reward.

The answer is no, and this is the biggest limit on the whole project. It is written
in my threat model file and I should have led with it.

In 28 of the 33 families neither arm learns anything. Both arms are policies I
wrote by hand. I wrote a script that behaves the way a cheating agent would behave,
and a script that behaves the way an honest agent would behave, and I ran both.
There is no learner at all. So no, they are not two seeds of the same learner. The
seeds change the layout and the noise, not any optimisation.

The other 5 families are different. They run real policy gradient against a reward
model that was actually fitted from preference data. There the cheating is not
scripted, it comes out on its own. But the hole in the reward model, and the label
saying which arm is which, are still mine.

So what does the benchmark show? It shows what you could detect \emph{in principle}
from each kind of information. It does not show what real trained agents leave
behind. Those are two different claims and I ran them together in earlier writeups.

What is the same in both arms: the task, the proxy reward function, and, after a
tuning step, the amount of proxy reward earned. What is different: the behaviour,
and the true reward.

\subhead{Why I force the visible scores to match}
If the cheating arm simply earned more visible score, a detector could read the
visible score and be right every time. It would have learned nothing about
cheating. It would have learned that my benchmark pays cheaters more. By making the
visible score the same, any detector that works has to have read something else.
And because I hand out the extra information in named channels, I know exactly
which one it read.

\heading{The four things a detector is allowed to see}
\label{h:levels}

You said each of these could mean a variety of things. Here is exactly what is in
each one.

\vspace{0.3em}
\begin{tabular}{@{}llp{7.7cm}@{}}
\toprule
Level & Field it adds & What the array actually holds \\
\midrule
L0 & \code{proxy\_rewards} &
  One number per episode: the proxy reward that episode. \\
L1 & \code{state\_counts} &
  Episodes $\times$ bins: how many times each state was visited that episode. \\
L2 & \code{behav\_trace} &
  One number per episode: a summary of the behaviour that episode, written by
  hand. On a gridworld it is (mean row visited) minus (mean column visited). \\
L3 & \code{true\_rewards} &
  One number per episode: the true reward. \\
\bottomrule
\end{tabular}
\vspace{0.3em}

Three things matter more than the table.

\subhead{The code deletes what you are not allowed to see}
Before a run reaches a detector it goes through a function that builds a fresh copy
with the higher fields emptied out. A detector at L0 cannot read the behaviour
number because there is nothing there to read. It is not a rule someone could
forget to follow.

\subhead{L3 is not a real level}
No real monitoring system gets the true reward. If you had it you would not need a
detector. L3 is there as a ceiling, to say how well detection could possibly do on
these tasks. It should never be quoted as a practical result.

\subhead{L2 is a number I chose, and that is a problem I have only half fixed}
I picked the behaviour number knowing which arm was the cheating one, and I picked
an axis I knew would separate them. The next two sections are about what that cost
me. Building L2 out of summaries that nobody chose with the answer in hand is work
I have not done yet.

\heading{What a detector is, and what fit means}
\label{h:detector}

A detector is a class with two methods.

\code{classify(run)} takes one run, already cut down to that detector's level, and
returns a number between 0 and 1. Higher means ``I think this one is cheating''. It
does not have to be a real probability. Only the ordering is used.

\code{detect\_onset(run)} returns which episode it thinks the cheating started, or
$-1$. This half of the project is not currently supported and I am not making
claims about it.

Most detectors are a fixed formula. Does the reward curve go flat and stay flat? Is
the reward spread lopsided? How far apart are the two state histograms? Nothing is
trained. The same formula is applied to every run.

\subhead{What fit means}
Some detectors also have a method called \code{fit}. It gets handed the two groups
with the labels attached, these are the cheating runs and those are the honest
ones, and it trains on them.

The benchmark decides how to score a detector with one line of code: does this
object have a method called \code{fit}? If not, run it on each run directly. If
yes, do 5-fold cross validation. Train on four fifths of the runs, score on the
fifth it did not see, then rotate.

Cross validation keeps this fair. Such a detector is never scored on runs it
trained on, so the number is real. But it answers a different question:

\begin{tabular}{@{}ll@{}}
\toprule
A fixed-formula detector & Does this formula separate the two arms? \\
A detector fitted on labels & Can the difference be learned from labelled examples? \\
\bottomrule
\end{tabular}

If I average both into one number per level, I answer neither. I report whichever
question happened to give the bigger number. That turns out to matter a lot, and it
is the section after next.

\heading{What the score means}
\label{h:score}

The score is AUROC. In plain words: pick one cheating run and one honest run at
random. AUROC is the chance that the detector gave the cheating one the higher
number. A tie counts as half. So 0.5 is a coin flip, 1.0 is always right, and 0.0
is always exactly wrong.

You asked whether a high score means a lot of reward hacking or a little. Neither,
and it is a fair thing to ask because the number looks like it should mean that.
AUROC measures how well a detector separates two groups. It says nothing about how
often hacking happens anywhere else. In my setup the mix is exactly half cheating
and half honest by construction, so the benchmark cannot say anything at all about
how common hacking is.

Scores are worked out per cell and then averaged upwards, per detector, per family,
per level.

\subhead{Where the 0.994 came from}
It was not the ordinary score. It was the transfer score. Fit the detectors on some
families, then test them on families they had never seen. If that number were real
it would mean something worth having: a detector trained on one kind of cheat
catches kinds it was never shown. Detection would carry across mechanisms. That was
the interesting claim.

It was not real.

\heading{Why the 0.994 was fake}
\label{h:bug}

You asked what I meant by ``a coding convention'', and whether I mean I set myself
a rule that changed what the code does. Yes. That is what happened, and you are
right that it is a bad thing to have done.

My contributor guide, which is a document about house style, said every family had
to produce its behaviour number with the same orientation: positive for the
cheating arm, negative for the honest arm. The reason was tidiness. It made
families comparable and plots consistent.

The main L2 detector scored a run by taking the mean of the last 100 entries of
that number. The benchmark worked out AUROC with cheating as the positive class.

Put those three facts together. The sign of the number was the label. The detector
was not reading behaviour. It was reading my style guide. A detector that
``transfers to a hacking mechanism it has never seen'' had worked out that all 33
families were written by one person following one rule.

\subhead{The sign I ignored for months}
Every family has a difficulty knob that sets how far apart I make the two arms. L2
AUROC sat at exactly 1.000 at every difficulty setting, in 12 of the 13 families
whose knob I swept. I moved the knob and the score did not move.

That cannot happen for a real measurement, and it is obvious once you say it out
loud. A knob that changes how far apart two things are cannot fail to move a
statistic that measures how far apart they are. Unless that statistic is measuring
something else. It was. The sign was already the answer, so distance did not
matter.

\subhead{The fix}
For each family and layout, a fixed rule gives either $+1$ or $-1$. Both arms get
multiplied by that one value before any detector sees them.

Multiplying \emph{both} arms is the whole trick. It flips the axis the pair is
defined on. It is not noise added to individual runs. If I flipped the two arms
separately I would not be hiding the direction, I would be destroying the pair,
because the opposite signs are what makes it a pair in the first place. Every size
is kept exactly, so the entry checks and the difficulty knob are untouched. The
only thing removed is the constant that let a detector read the label off a sign.

The transfer score falls from 0.994 to 0.508. The L2 average falls from 0.7229 to
0.5629.

\heading{Am I sure about the 0.508?}
\label{h:sure}

You asked this straight, so here is what I have and what I do not.

\subhead{What backs it up}
The comparison is like for like. The randomisation is a single flag. Turn it off
and 0.994 comes back exactly. And 0.508 is not one run. I repeated the whole thing
20 times with different layout and rollout seeds, which is 600 detector passes, and
got 95\% intervals by resampling those 20 repeats.

\subhead{What chance means here}
0.508 against 0.5 is not something I can tell apart from a coin flip with this
sample. I am not claiming the detectors are slightly better than chance.

\subhead{What would overturn it}
The randomisation hides the direction inside any single run. But a detector is also
allowed to look at a whole batch of runs without labels, and it could in principle
work the direction out from the batch. That is allowed and I have not shown that
nobody can do it. A detector that did that and still transferred would put the
number back up, and it would deserve to.

\subhead{One hole I know about and have not closed}
At L2 the detector also sees the state counts, and I do not flip those. A count
over one family's own grid has no direction to flip. So a family whose bins line up
with the same axis as its behaviour number still carries, in principle, enough to
recover the old orientation. Nothing in my repo does this. Doing it on purpose
would be reading the label by a longer route. It only closes properly when I
rebuild L2.

\subhead{What the fix does not do}
It removes the free direction and nothing else. The behaviour number is still a
feature I picked knowing the answer. So an L2 score under randomisation means
``this detector did not need the sign rule''. It does not mean ``L2 is blind to the
label''.

\heading{Five of the thirty}
\label{h:five}

Five of the thirty \emph{detectors}, not the environments. They are:

\vspace{0.3em}
\begin{tabular}{@{}ll@{}}
\toprule
Detector & Level \\
\midrule
Reward MLP & L0 \\
State Coverage Rate & L1 \\
State Divergence & L1 \\
Visitation Entropy Trend & L1 \\
Trajectory MLP & L2 \\
\bottomrule
\end{tabular}
\vspace{0.3em}

These are the ones with a \code{fit} method, so they are trained on labels and
scored by cross validation.

This matters because of what my main table was being read as. The table is a
ladder. Give a detector more information, does it do better? The answer looked like
yes, with a big jump from reward-only to state counts.

Split the two kinds of detector apart, on the same 20 repeats:

\vspace{0.3em}
\begin{tabular}{@{}lcc@{}}
\toprule
Level & Best of all detectors & Best of the fixed-formula ones \\
\midrule
L0 (reward only)      & 0.5477 & 0.5437 \\
L1 (+ state counts)   & 0.9416 & 0.5340 \\
L2 (+ behaviour)      & 0.9430 & 0.5872 \\
L3 (+ true reward)    & 0.9746 & 0.9746 \\
\bottomrule
\end{tabular}
\vspace{0.3em}

And each step of the ladder, measured across the 20 repeats:

\vspace{0.3em}
\begin{tabular}{@{}llcc@{}}
\toprule
Step & Detectors & Gap [95\% interval] & Repeats where the step held \\
\midrule
L1 $-$ L0 & all            & $+0.394$ \ $[+0.384,\ +0.403]$ & 20 of 20 \\
L1 $-$ L0 & fixed-formula  & $-0.010$ \ $[-0.028,\ +0.008]$ & 5 of 20 \\
\addlinespace
L2 $-$ L1 & all            & $+0.001$ \ $[-0.004,\ +0.006]$ & 13 of 20 \\
L3 $-$ L2 & all            & $+0.032$ \ $[+0.027,\ +0.037]$ & 20 of 20 \\
\bottomrule
\end{tabular}
\vspace{0.3em}

Read the first two rows together. The gap flips sign and stops being a gap. The
whole published climb from reward-only to state counts was one detector that had
been trained on labels. That is a real result about my benchmark. It is not the
result the ladder was being read as, and nobody could have got it out of a
single-column table.

The third row matters too. The step from state counts to behaviour has an interval
that includes zero, so there is no measurable step there either.

\heading{What I meant by admission criteria}
\label{h:checks}

You said the phrase sounded weird and had probably come from an LLM. Fair. Here is
the idea without it.

A pair is a claim. I am saying these two arms earn the same visible score while
doing different things. A claim is not a fact. So before a pair goes into the set,
I test it. Six tests per cell:

\begin{enumerate}
\item The two arms earn the same visible score on average.
\item Their visible scores have the same \emph{shape}, not just the same average.
      Checked three ways.
\item The behaviour really is different.
\item The true rewards really are different.
\item You can find the moment the cheating starts, and find it consistently.
\item The cheating keeps going to the end instead of fading out.
\end{enumerate}

\subhead{What test 1 used to be}
It was literally this: \code{abs(mean\_auroc - 0.5) < 0.10}.

That is: run a reward-only detector, and if it fails to tell the arms apart, decide
the arms match. Which is ``I did not find a leak'' written down as ``there is no
leak''.

That is only as good as the experiment's ability to find a leak, and mine had
almost none. At four seeds a side the standard error is 0.0625. So:

\vspace{0.3em}
\begin{tabular}{@{}lc@{}}
\toprule
If the true value were\dots & \dots the old test passed it this often \\
\midrule
0.500 (really matched)   & 89.6\% \\
0.556                    & 76.2\% \\
0.611 (\textbf{leaking}) & \textbf{43.5\%} \\
0.714                    & 2.5\% \\
\bottomrule
\end{tabular}
\vspace{0.3em}

A family that was clearly leaking passed nearly half the time. Worse, the incentive
runs backwards. Noisier data made the test \emph{easier} to pass, because a noisy
measurement lands near 0.5 more often.

\subhead{The fix, and what it cost}
Swap the test round.

\vspace{0.3em}
\begin{tabular}{@{}ll@{}}
\toprule
The old test asked & Can I show these two are different? \\
The new test asks  & Can I show these two are close? \\
\bottomrule
\end{tabular}
\vspace{0.3em}

Now the whole confidence interval has to fit inside 0.5 plus or minus 0.10. Noise
makes a wide interval, and a wide interval does not fit. So noise now makes it
harder to pass, which is the right way round.

The margin stayed at 0.10, so the claim itself did not change. What changed is how
much data it takes, and that change is big:

\vspace{0.3em}
\begin{tabular}{@{}lcc@{}}
\toprule
& Before & After \\
\midrule
Seeds per side per layout & 4 & 24 \\
Layouts & 12 & 12 \\
Runs per cell & 96 & 576 \\
Half width of the interval & 0.117 & 0.046 \\
\bottomrule
\end{tabular}
\vspace{0.3em}

Look at that 0.117 against the margin. The interval was wider than the whole 0.10
band. So the old setup could not have passed the new test at any measured value,
not even at exactly 0.5. That is the numeric reason a weak test was standing in for
the right one. The right test would have failed everything.

\heading{The check that broke inside the fix}
\label{h:degen}

This is the part I would show a student, because it is the same mistake happening
again one level down, inside its own repair.

AUROC only looks at order. If every score it has to rank is tied, the standard
function returns exactly 0.5, because ties count half. Not because the two arms
were compared and found equal. Because nothing was compared.

Now feed that into the new test. Every layout gives exactly 0.5. So the spread is
0. So the interval collapses to the single point $[0.5000, 0.5000]$. So it fits
inside any margin at all, 0.10 or 0.001 or anything smaller. The test passes
always, and can never fail.

That is exactly the problem the new test was brought in to fix. The first ledger
the new gate produced passed 15 of its 35 cells this way. Every one of them was a
family whose proxy reward is a constant by design. A constant score is trivially
matched. It is also useless, and the gate could not tell those two apart.

\subhead{Two guards, and a third answer}
So before reading any of these tests, the gate now asks two questions.

\emph{Could the number have come out different?} For each layout it works out what
fraction of the cross-arm run pairs the detector actually puts in an order. If the
range the score could possibly reach already fits inside the band, the test was
decided before the first run.

\emph{Did the signal mean anything?} The first guard is cheap to fool. A proxy
reward of $0.675$ plus a wobble of size $0.0000001$, which is a constant plus noise
a hundred million times smaller than the reward and which nothing could act on,
pushes the first guard to 0.961 and the whole gate says pass. No tie tolerance
fixes that, because a wobble ten times bigger than the tolerance always clears it.
So the second guard is about the family and not the scores: how big the proxy
reward's spread is compared to the reward itself, with a floor underneath it.

Neither guard covers the other. On the current grid each one catches three families
the other cannot see.

And there are now three answers instead of two:

\vspace{0.3em}
\begin{tabular}{@{}lp{9.2cm}@{}}
\toprule
Pass & Measured, and matched inside the margin. \\
Fail & Measured, and outside the margin. The score leaks, and every other
       number on that cell is suspect. \\
Could not measure & The test could not be run at all. Says nothing, in either
       direction. \\
\bottomrule
\end{tabular}
\vspace{0.3em}

``Could not measure'' is not a soft fail and it is not a pass. Filing it as a pass
is the mistake this whole audit is about, and I made it twice at two different
levels. Before the fix: ``I did not find a leak'' published as ``there is no
leak''. After the fix: ``I could not look'' published as ``I looked and it was
fine''. Same mistake. The second one survived a repair aimed at the first.

\subhead{Where it stands now}
Of the 50 cells in scope: \textbf{15 pass, 30 could not be measured, 5 genuinely
failed}. Passing is the minority. The count dropped from an earlier ``30 of 35
passed'' not because any family got worse, but because the new guard stopped
counting cells whose tests were never really run.

And the scope is small, which I would rather tell you than have you find. The
ledger covers 10 of the 33 families. The other 23 are unchecked, which is not the
same as checked and fine.

\heading{What I actually want to work on}
\label{h:point3}

You asked whether my interest is in building benchmarks or in experimental rigour,
how I would change things, and, the sharpest one, whether I am sure there is a real
problem in the literature.

Let me take the last one first. No, I am not.

The paper claims a kind of failure: checks that cannot come back negative, given
how the benchmark is built. It gives three ways to find them. It also has a survey
of 19 published benchmarks, and I wrote down what I was going to look for before I
started so I could not move the target afterwards.

I did not find it. The pattern did not turn up in the other benchmarks in the form
I found it in mine. That is in the paper as a claim I withdrew.

So my honest position is this. I found this in one benchmark, and it was mine. I
went looking for it in nineteen others under rules I fixed in advance, and it was
not there. What is left is the method and one detailed case study. Not a claim
about the field.

A second prediction of mine failed recently, and I am telling you for the same
reason. Earlier work of mine said the ``true rewards differ'' check is blind to
cheating that \emph{adds to} a competent agent rather than \emph{replacing}
competence. I predicted that two categories in an outside dataset would be the
adding kind. They are not. All of them avoid doing the real work. The prediction
was wrong.

Here is the thing I do want to work on, stated concretely enough that you can
attack it:

\vspace{0.4em}
\begin{quote}
\emph{When an evaluation gives you a number, is that number evidence about the
thing the evaluation says it measures?}
\end{quote}
\vspace{0.4em}

That splits into two things I can actually test.

First, before reporting any score, can you show that a benchmark's entry checks are
capable of rejecting something? A criterion is not a check until it comes with a
demonstration that it can fail. A control is decoration unless the failures get
published next to the passes.

Second, does a result that holds on the data you have survive a change you named in
advance? The 0.994 to 0.508 collapse is one case of that. I named the change first,
and the result did not survive it.

You are right that people already check their own tests, more or less carefully. My
claim is smaller than ``nobody does this''. It is that auditing your own benchmark
is a one-off thing you write a paper about, rather than a routine step with a
standard shape. I do not have a framework for making it routine. I have one worked
example and three checks.

On your hyperparameter point, you are right, and it made me drop a word I had been
using loosely. Nothing survives an arbitrary change, and any algorithm can be made
badly behaved on purpose. So ``is this robust'' has no answer. ``Does this result
survive this specific change, named in advance'' does.

\heading{What this project does not show}
\label{h:limits}

All in one place, so you do not have to trust me to bring them up one at a time.

\begin{itemize}
\item \textbf{The cheating is scripted.} 28 of 33 families are policies I wrote by
      hand. This is about what a channel holds in principle, not about real trained
      agents.
\item \textbf{The public table is one draw.} Every cell is a single 5-against-5
      comparison, run once. The standard error near 0.5 is about 0.19. Do not read
      the third decimal, and do not rank two detectors that differ by less than
      about 0.2. This is fixed for the summary numbers, which come from the 20
      repeats, but not for individual cells.
\item \textbf{All 30 detectors share that one draw}, so their cell values move
      together and their differences are not independent.
\item \textbf{The set I check is not the set I score.} The entry checks average
      over 12 layouts. The public table uses one.
\item \textbf{L0 sitting at chance is a control, not a discovery.} The families are
      \emph{built} so reward-only detectors cannot work. An earlier version of my
      README reported this as a finding that ``reward-only monitoring fails''. That
      was wrong. It is only useful in the failure direction: an L0 detector going
      above chance means the construction leaked.
\item \textbf{One row in the table is a duplicate.} The detector I labelled
      ``Perfect Feature Oracle'' at L3 is a subclass of an L2 detector that
      overrides only its name and its level. It never reads the true reward. It
      agrees with its parent on 33 of 33 families. I keep it as a labelled
      cross-check and leave it out of every average, but the old ``L3 mean over
      two detectors'' figure averaged one oracle with one copy of an L2 baseline
      and measured nothing.
\item \textbf{The onset half is unsupported.} I published this as ``Reward Hacking
      \emph{Onset} Benchmark''. That column is not currently backed by anything. I
      keep the short name so old links work, and the long name should not be used.
\item \textbf{The entry checks are barely covered by my own tests.} The full check
      is never run on a real registered family anywhere in the test suite. The two
      tests that touch it use made-up fixtures. A green test suite here means ``the
      screens that exist passed'', not ``the families are checked''.
\end{itemize}

\heading{Your questions, and where I answered them}
\label{h:index}

\begin{tabular}{@{}p{10.4cm}r@{}}
\toprule
Your question & Page \\
\midrule
What is a cheating run? What is an honest run? & \pageref{h:words} \\
How do you distinguish between them? & \pageref{h:learn} \\
Are the learning algorithms the same? The reward functions? & \pageref{h:learn} \\
Can these be seen as two seeds optimising the same reward? & \pageref{h:learn} \\
What is a proxy reward function? & \pageref{h:words} \\
What is a finished run? Policy, trajectory, or distribution? & \pageref{h:words} \\
What are reward-only, state visitation, behavioural, oracle? & \pageref{h:levels} \\
Cheating is contested in the literature. How can you detect it? & \pageref{h:learn} \\
What is a detector? & \pageref{h:detector} \\
What does it mean to fit a detector? & \pageref{h:detector} \\
What are these scores? Is 0.994 an average of them? & \pageref{h:score} \\
Would a real 0.994 mean a lot of hacking, or a little? & \pageref{h:score} \\
What do you mean by a coding convention? Was it a bug? & \pageref{h:bug} \\
A style rule that changed what the code does? & \pageref{h:bug} \\
What is the significance of 0.508? Are you sure it is correct? & \pageref{h:sure} \\
What does confidence from twenty replications mean? & \pageref{h:sure} \\
Five of the thirty. Detectors or environments? & \pageref{h:five} \\
What is an admission criterion? & \pageref{h:checks} \\
Is your interest benchmarks, or experimental rigour? & \pageref{h:point3} \\
How would you change the status quo? What would the impact be? & \pageref{h:point3} \\
Are you sure there is a real problem in the literature? & \pageref{h:point3} \\
\bottomrule
\end{tabular}

\heading{If you want to check any of this}
\label{h:sources}

Everything above is in the repo at \url{https://github.com/Aarav500/rhob}. The
files that matter:

\begin{itemize}
\item \code{README.md} for the correction notice, the method, and the limits.
\item \code{docs/THREAT\_MODEL.md} for what the benchmark does and does not
      measure, including the scripted-policy limit.
\item \code{src/rhob/v3/sign\_randomization.py} for the fake 0.994, written up at
      the top of the file.
\item \code{src/rhob/v3/admission\_gate.py} for the entry checks.
\item \code{src/rhob/v3/access.py} and \code{src/rhob/detectors/posthoc.py} for the
      four levels and the run format.
\item \code{src/rhob/detectors/supervision.py} for the five fitted detectors.
\item \code{leaderboard/v5\_replicated.json} for the 20 repeats and the intervals.
\item \code{admission/ADMISSION\_LEDGER.md} for every cell and every check, the
      failures alongside the passes.
\end{itemize}

\end{document}
